\section{Private serving and private delivery}
\label{sec:tech:pir}

This section covers the half of the system that is built and running: the user
is scored against the catalogue without any server learning the user's
embedding, and the resulting records are fetched without either serving party
learning which records were fetched. Private \emph{training} is the next phase
and is not claimed here.

\subsection{The catalogue is fixed-width, and that is a security property}
\label{sec:tech:pir:catalogue}

A private read must return the same number of bytes whatever it reads, or the
response length itself leaks the record. MovieLens-100K's \texttt{u.item} is
variable-length text, so it is packed into a fixed record of $L = 256$ bytes.

Two fields are dropped. Field~3, the video release date, is \emph{empty on all
1682 lines} --- measured, not assumed --- and field~4, the IMDb URL, is the
largest field and is used by nothing downstream. What remains is
\texttt{id\,|\,title\,|\,date\,|\,19 genre flags}, whose longest instance is
117 bytes of content, leaving 139 bytes of headroom inside $L$. Every record is
checked against that budget at load time, so a future dataset with a longer
title fails loudly at startup rather than truncating silently at answer time.

Fields are carried with explicit one-byte length prefixes rather than a
delimiter. Nine titles contain Latin-1 bytes and the file does not decode as
UTF-8 at all, so it is read as bytes and never re-encoded; with arbitrary bytes
in the payload, any delimiter choice risks colliding with content, whereas a
length prefix cannot. Three records break a naive parser and each is handled
explicitly: id~267 is \texttt{unknown} with an empty date and empty URL,
id~1242 is the width driver, and one release date is \texttt{4-Feb-1971} rather
than the \texttt{DD-Mon-YYYY} that every other date follows --- which never
mattered, because dates are carried as opaque text and never parsed into a
structured type.

The 1682 items pad to a domain of $2048$, so $\mathrm{db} = 11$. Padding
records are all-zero. This leaks nothing precisely because record width is
fixed regardless of content: a fetch of a padding slot is indistinguishable
from a fetch of a real one.

\subsection{The private read}
\label{sec:tech:pir:read}

The client wants record $\alpha$ without either server learning $\alpha$. It
generates a DPF key pair for the point function that is $1$ at $\alpha$, sends
one key to each of $\party{0}$ and $\party{1}$, and each server expands its key
over the whole domain and takes an inner product against the catalogue,

\[
  \mathrm{ans}_c[w] \;=\; \sum_{j=0}^{2^{\mathrm{db}}-1}
    \EvalFull(k_c)[j] \cdot D[j][w],
  \qquad c \in \{0,1\},
\]

over the $W = L / (\bits/8)$ ring words of a record ($W = 32$ at $\bits = 64$).
Because $\EvalFull(k_0)[j] - \EvalFull(k_1)[j]$ is $1$ at $j = \alpha$ and zero
elsewhere, the \emph{difference} of the two answers is exactly $D[\alpha]$.
Note the difference rather than the sum: this project omits the $(-1)^c$ factor,
per Section~\ref{sec:tech:dpf:sign}.

Two properties carry the obliviousness claim, and they are worth separating
because the second is the stronger one:

\begin{enumerate}
  \item The server touches \emph{every} record, in index order, with no
    data-dependent branch. The loop bound is the domain size and nothing inside
    the loop depends on $\alpha$.
  \item Each expanded coefficient is a \emph{share}, so it is pseudorandom at
    every index rather than zero at all but one. There is no data pattern for a
    compiler to exploit or a cache-timing observer to see, even in principle.
    This is strictly more than ``the code contains no branch''.
\end{enumerate}

The system has three servers but the DPF is $(2,2)$. We fix $\party{0}$ and
$\party{1}$ as the key holders; all three hold the catalogue, which is public
data, so replication costs nothing in privacy.

\paragraph{What it costs.} At $\mathrm{db} = 11$ a query is a single
\textbf{227-byte} key per server against a 2048-entry domain
(Table~\ref{tab:tech:keysize}), against 16~KB for the naive one-hot query the
same read would otherwise need: a $72\times$ reduction in upload. The server's
work is one $\EvalFull$ plus $2^{11} \times 32$ ring multiply-accumulates.

\paragraph{Verification.} Reads are checked byte-exactly against a cleartext
lookup: 60 random indices at $\bits = 64$ and 20 at $\bits = 128$, plus both
domain endpoints, padding slots, and rejection of a key whose domain width does
not match the server's.

One test in this group was wrong, and establishing that took more work than
accepting it would have. A natural check --- ``no word of one server's share
should equal the plaintext'' --- fails, with 18 of 32 words matching. That is
not a leak. Records pad to 256 bytes while content never exceeds 117, and the
trailing genre bytes are usually zero, so the upper words are identically zero
across the \emph{entire} catalogue; any linear combination of zeros is zero,
for every $\alpha$ equally. It is public structure, not a signal. The test now
derives from the data which words actually vary --- 14 of 32 at $\bits = 64$
and 7 of 16 at $\bits = 128$, which agree at 112 bytes and so cross-check each
other --- and asserts only on those.

\subsection{Scoring under secret sharing}
\label{sec:tech:pir:serving}

User embeddings $\userembed$ are secret; item embeddings $\itemembed$ are
public, which is \textsc{Nudge}'s design and is what makes its power iteration
cheap. The user's row is shared 2-of-3, and because $\itemembed$ is public,
$\mathbf{s} = \avec{i} \cdot \itemembed$ is a \emph{linear map applied to a
shared vector by a public matrix}. Each server computes its own share of all
$\items$ scores entirely locally: no communication, no correlated randomness,
no multiplication protocol, no rounds. This is the single largest saving in the
serving path and it comes directly from $\itemembed$ being public.

\paragraph{Scores stay at $2\fracbits$ scale.} Both factors are encoded at
$\fracbits$ fractional bits, so the product carries $2\fracbits$. Truncating
back to $\fracbits$ requires $\mathrm{Trunc}_\fracbits$, which is interactive
and belongs to the training phase, so serving decodes at $2\fracbits$ instead.
This is recorded rather than left implicit because getting it wrong scales
every score by $2^{20}$, which is \emph{monotone} and therefore completely
invisible in a ranking.

\paragraph{The seen-item mask.} The ideal functionality sets already-rated
items to $-\infty$, which a finite ring cannot represent. We use an additive
shared vector carrying $-(2^{55})$ at seen positions and zero elsewhere;
servers add it to their score shares, which is again local. Observed encoded
scores peak near $9.7 \times 10^{12}$ and the ring floor is near
$-9.2 \times 10^{18}$, so the sentinel sits about $3700\times$ below any real
score and about $256\times$ above the floor. Both margins are asserted by a
test rather than argued in prose. The comparison in the top-$\topk$ selection
is \emph{signed}: comparing these two's-complement values as unsigned would
sort masked items \emph{first}, and a smoke test inspecting only the head of
the list would not notice.

We state one thing plainly rather than overclaim it: in this phase the user
reconstructs the score vector anyway, so masking server-side is
\emph{equivalent} to masking client-side. It is implemented server-side because
that is what the ideal functionality specifies and what the composed system
will need once the servers do more than a linear map --- not because this phase
requires it.

\subsection{The float-to-ring boundary}
\label{sec:tech:pir:boundary}

Exactly one program converts a float into a ring element:
\texttt{model/export.py}. Everything downstream trusts it, so it is tested as a
contract rather than assumed. A committed file of 25 \texttt{<double> <int64>}
pairs --- zero, both signs, the quantum and half the quantum, the observed
extremes of $\userembed$, $\itemembed$ and the scores, and the two's-complement
edge --- is read by the C++ test suite, so the boundary is checked on a fresh
clone with no export run.

It found two bugs immediately, both of which would have been hard to find
later.

\textbf{Rounding mode.} \texttt{numpy.rint} and Python's \texttt{round} use
banker's rounding, half to even, so $\mathrm{rint}(0.5) = 0$; C++
\texttt{std::round} rounds half away from zero, so $\mathrm{round}(0.5) = 1$.
They disagree at every exact half-quantum. The contract test caught this on its
first run. Python was changed to match C++, since C++ is what ships.

\textbf{Quantise-then-multiply is not multiply-then-quantise.} The reference
scores were originally computed by encoding the float product; the protocol
multiplies encoded factors, which rounds $\dime$ times and then accumulates
exactly. The two disagreed on \textbf{1650 of 1682 scores}, by up to
$2.1 \times 10^{7}$ out of $4.3 \times 10^{12}$. That is about $5\times10^{-6}$
relative and irrelevant to the ranking, which is exactly why it matters: any
tolerance-based comparison would have accepted it, and the oracle would have
been quietly wrong for every later exact test. The reference now computes what
the protocol computes, in arbitrary-precision integers.

Serving is therefore checked against the oracle by \textbf{exact} equality, not
by tolerance. A linear map over a ring has no tolerance to spend, and a
tolerance here would hide precisely the endianness, sign-extension and
scale bugs the boundary is prone to. All 1682 scores match bit-for-bit.

\paragraph{Ring width, measured.} An open question in the design was whether
$\bits = 64$ suffices at MovieLens scale or whether $\bits = 128$ is needed.
Measured on the actual factorisation: $\userembed$ peaks at $61.4$,
$\itemembed$ at $0.219$ (every row has unit $L_2$ norm by construction), and at
$\fracbits = 20$ the worst-case $\dime = 16$ accumulation is
$\mathbf{47.8}$ \textbf{bits of the 63 available} --- roughly a $32{,}000\times$
margin. Quantisation costs $\max|S_q - S| = 5.3 \times 10^{-5}$, far below the
gaps between top-20 scores, so the ranking is preserved. So $\bits = 64$ is
sufficient at this scale, and the implementation remains templated on the ring
so the $\bits = 128$ arm is tested alongside it at every step.

\subsection{End to end}
\label{sec:tech:pir:demo}

\texttt{demo --user 42 --k 10} runs the whole path: the embedding is split
across three servers, each scores the catalogue locally, the user reconstructs
and selects the top ten, and each of the ten records is then fetched by a
two-server DPF-PIR read. It prints ten real film titles, and it checks itself
twice: every fetched record is compared byte-for-byte against a cleartext
lookup, and the ranking is compared position-for-position against the
independently computed Python reference. Both agree.

For user~42, masking the 112 items already rated in the training split, the top
ten are \emph{The Princess Bride}, \emph{Indiana Jones and the Last Crusade},
\emph{Toy Story}, \emph{Jerry Maguire}, \emph{Apt Pupil}, \emph{Good Will
Hunting}, \emph{Back to the Future}, \emph{The Rock}, \emph{Mission:
Impossible} and \emph{Aladdin} --- fetched at a cost of one 227-byte key per
server per title, with neither server able to say which titles they were.
